% !TeX root = ../初等数学笔记.tex
% 由主文档通过 \input 载入；请编译主文档。
\setcounter{chapter}{19} % 保留原书章号。
\chapter{圆锥曲线}
\label{chap:em-conics}
曲线方程应当是几何条件的等价表达。距离定义往往含有根式，推导方程需要平方，但反向检验必须恢复被平方隐藏的符号和位置限制。

\section{从几何定义到曲线方程}
\emph{圆锥曲线}（conic section）包括圆、椭圆、双曲线和抛物线。本章从焦点距离及焦点准线关系出发讨论非退化实曲线；任意二元二次方程的完整分类需要进一步的代数工具。

求轨迹的一般步骤是选择坐标、表示距离、保留定义域和位置条件、消去根式、验证所得方程上的每一点满足原条件。若原题排除了某条线或某个点，消元后仍须排除。坐标轴应尽量沿对称轴设置，使对称性体现在方程中。

\section{椭圆的定义与标准方程}
给定焦点 $F_1=(-c,0),F_2=(c,0)$，$c>0$。到两焦点距离之和为 $2a$ 的点集称为\emph{椭圆}（ellipse），要求 $a>c$。设两距离为 $r_1,r_2$，由
\[
 r_1+r_2=2a,\qquad r_1^2-r_2^2=4cx
\]
得到 $r_1=a+cx/a,\ r_2=a-cx/a$。代入距离平方并整理，得到
\begin{equation}\label{eq:em-ellipse}
 \frac{x^2}{a^2}+\frac{y^2}{b^2}=1,\qquad b^2=a^2-c^2>0.
\end{equation}
反之，满足该式的点有 $|x|\le a$，故 $a\pm cx/a\ge a-c>0$；它们的平方分别等于两距离平方，所以确为两距离，和为 $2a$。这样完成了必要性与充分性。

椭圆关于两坐标轴和原点对称，顶点为 $(\pm a,0),(0,\pm b)$，长、短轴长为 $2a,2b$。离心率 $e=c/a\in(0,1)$ 描述扁平程度，焦点重合的 $c=0$ 特例成为圆。以右焦点为例，$PF_2=a-cx/a=e(a^2/c-x)$，后者等于 $e$ 倍的点到准线 $x=a^2/c$ 的距离，从而也得到焦点准线表述。

\section{双曲线的定义、两支与渐近线}
同样取两焦点，距离差绝对值为 $2a$ 的点集称为\emph{双曲线}（hyperbola），其中 $0<a<c$。对右支先写 $r_1-r_2=2a$，由平方差得 $r_1+r_2=2cx/a$，所以 $r_1=cx/a+a,r_2=cx/a-a$，推导
\begin{equation}\label{eq:em-hyperbola}
 \frac{x^2}{a^2}-\frac{y^2}{b^2}=1,\qquad b^2=c^2-a^2.
\end{equation}
右支满足 $x\ge a$，因此候选距离均为正，反向平方检验成立；左支由关于纵轴的对称获得。方程自动给出 $|x|\ge a$，两顶点为 $(\pm a,0)$，离心率 $e=c/a>1$。

渐近线为 $y=\pm(b/a)x$。例如右上支 $y=(b/a)\sqrt{x^2-a^2}$，在 $x\ge a$ 时
\[
 \frac ba x-y=\frac{ab}{x+\sqrt{x^2-a^2}}\longrightarrow0
 \quad(x\longrightarrow+\infty).
\]
因而点到对应直线的距离也趋于零；这给出了“靠近渐近线”的准确含义，不表示曲线最终与直线重合。

\section{抛物线的定义与参数约定}
到定点与不经过它的定直线距离相等的点集称为\emph{抛物线}（parabola）。取焦点 $(p,0)$、准线 $x=-p$，其中 $p>0$，距离关系为
\[
 \sqrt{(x-p)^2+y^2}=|x+p|,
\]
平方整理得到
\begin{equation}\label{eq:em-parabola}
 y^2=4px.
\end{equation}
反向看，方程推出 $x\ge0$，故 $x+p>0$，距离平方相等能恢复原等式，没有多余分支。顶点为原点，对称轴为横轴，开口向右。其他方向通过交换坐标或变号处理；本书参数始终采用 $4p$ 的系数，故焦点横坐标为 $p$，不能混用其他记号体系。

抛物线的焦点准线距离比恒为 $1$，对应离心率 $1$。写成 $x=y^2/(4p)$ 后，它也是以 $y$ 为输入的二次函数图像；是否能写成 $y=f(x)$ 则取决于是否选择上下分支。

\section{直线交点、切线与综合轨迹}
将参数直线 $P=P_0+t\boldsymbol v$（$\boldsymbol v\ne0$）代入曲线，若得到 $At^2+Bt+C=0$ 且 $A\ne0$，就由判别式判断两个不同交点、重合交点或无实交点。两不同交点间弦长为 $|\boldsymbol v|\sqrt{B^2-4AC}/|A|$，中点参数为 $-B/(2A)$。若 $A=0$，必须改作一次或常数方程讨论。

在非退化圆锥曲线上，相切的代数判定是交点对应二重根，而不能只数不同交点。例如横轴 $y=0$ 与 $y^2=4px$ 只有一个交点，却是对称轴；顶点切线是 $x=0$。通过点 $(x_0,y_0)$ 的切线公式分别为
\[
 \frac{x_0x}{a^2}+\frac{y_0y}{b^2}=1,\qquad
 \frac{x_0x}{a^2}-\frac{y_0y}{b^2}=1,\qquad
 y_0y=2p(x+x_0),
\]
依次对应式 \eqref{eq:em-ellipse}、\eqref{eq:em-hyperbola}、\eqref{eq:em-parabola}，且给定点必须满足相应曲线方程。这些式子可由直线代入后要求对应参数的常数项和一次项均为零、只剩二次项来验证，无需导数。双曲线中平行于渐近线的方向可能使二次项消失，单个交点不代表二重根。

\section{例题、习题与学习检查}
\begin{example}
椭圆 $x^2/9+y^2/4=1$ 的焦点与离心率是什么？求它与横线 $y=1$ 所截弦长。
\end{example}
\begin{solve}
$a=3,b=2,c=\sqrt5$，焦点为 $(\pm\sqrt5,0)$，离心率为 $\sqrt5/3$。代入 $y=1$ 得 $x^2/9=3/4$，两交点横坐标为 $\pm3\sqrt3/2$，弦长为 $3\sqrt3$。
\end{solve}
\begin{enumerate}
\item 求 $x^2/4-y^2/5=1$ 的焦点和渐近线。答：$(\pm3,0)$，$y=\pm(\sqrt5/2)x$。
\item 求 $y^2=8x$ 的焦点、准线及点 $(2,4)$ 处切线。答：$(2,0)$、$x=-2$、$y=x+2$。
\item 两焦点距离为 $6$，距离和为 $4$，能否形成椭圆？答：不能，距离和不小于焦点间距。
\end{enumerate}

